\section{Тест производительности}
Тест производительности сравнивает решение на динамическом программировании с наивным рекурсивным перебором без запоминания подзадач.
Наивный вариант считает стоимость по тому же соотношению, что и ДП: $f(n) = n + \min(f(n-1), f(n/2), f(n/3))$ с учётом делимости,
но каждый раз вычисляет $f$ заново. Поэтому одни и те же значения пересчитываются многократно, а количество рекурсивных вызовов
$T(n) = T(n-1) + T(n/2) + T(n/3) + 1$ растёт сверхполиномиально.

Программа собрана компилятором MSVC 19.51 в конфигурации Release (флаг \texttt{/O2}), стандарт C++20.
Время измерялось с помощью \texttt{std::chrono::steady\_clock}, каждый замер повторялся 5 раз,
в таблице указано среднее значение. Измерялось только вычисление минимальной стоимости, без чтения ввода и без вывода.

\begin{table}[h]
\centering
\begin{tabular}{|c|c|c|c|}
\hline
$n$ & ДП, с & Наивная рекурсия, с & Отношение \\
\hline
50  & 0{,}000000760 & 0{,}000005600 & 7{,}4 \\
75  & 0{,}000000520 & 0{,}000023280 & 44{,}8 \\
100 & 0{,}000000500 & 0{,}000076760 & 153{,}5 \\
500 & 0{,}000001580 & 0{,}254238580 & 160\,910{,}5 \\
750 & 0{,}000003880 & 3{,}056941520 & 787\,871{,}5 \\
\hline
\end{tabular}
\end{table}

\begin{table}[h]
\centering
\begin{tabular}{|r|c|c|c|}
\hline
$n$ & ДП, с & Рост к предыдущему & Время на элемент, нс \\
\hline
$10^3$ & 0{,}000008860 & --- & 8{,}86 \\
$10^4$ & 0{,}000040460 & 4{,}57 & 4{,}05 \\
$10^5$ & 0{,}000310240 & 7{,}67 & 3{,}10 \\
$10^6$ & 0{,}004571280 & 14{,}74 & 4{,}57 \\
$10^7$ & 0{,}048045160 & 10{,}51 & 4{,}80 \\
\hline
\end{tabular}
\end{table}

\begin{alltt}
laba-7>cmake -S . -B build
laba-7>cmake --build build --config Release
laba-7>build/Release/bench.exe
n       ДП, с           наивная, с   во сколько
50      0.000000760     0.000005600  7.4
75      0.000000520     0.000023280  44.8
100     0.000000500     0.000076760  153.5
500     0.000001580     0.254238580  160910.5
750     0.000003880     3.056941520  787871.5
n       ДП, с
1000    0.000008860
10000   0.000040460
100000  0.000310240
1000000 0.004571280
10000000 0.048045160
\end{alltt}

Из первой таблицы видно, что время наивной рекурсии растёт очень быстро. При увеличении $n$ с 50 до 75 оно возрастает примерно в 4 раза,
с 75 до 100~--- в 3{,}3 раза, при росте $n$ с 100 до 500 (в 5 раз)~--- более чем в 3000 раз, а при росте с 500 до 750
(в 1{,}5 раза)~--- ещё в 12 раз. При $n = 750$ перебор работает около трёх секунд, а при $n = 10^7$ он заведомо неприменим.
Такой рост, намного более быстрый, чем у любого степенного закона, подтверждает, что без запоминания подзадач задача решается за
сверхполиномиальное время. Отношение времени перебора ко времени ДП вырастает с 7{,}4 раза при $n = 50$ до почти $8 \cdot 10^5$ раз
при $n = 750$, поэтому запоминание подзадач необходимо.

Во второй таблице время ДП измерено для $n$ от $10^3$ до $10^7$. На малых $n$ (до 100) время ДП составляет доли микросекунды и
сравнимо с погрешностью таймера и накладными расходами на выделение памяти под вектор, поэтому небольшие колебания
(например, время для $n = 75$ меньше, чем для $n = 50$) не отражают реальной зависимости. На больших $n$ картина устойчива: при каждом
увеличении $n$ в 10 раз время растёт в 4--15 раз, а время, приходящееся на один элемент таблицы, остаётся в пределах 3--5~нс.
Это согласуется с оценкой $O(n)$: каждая ячейка вычисляется за $O(1)$.

Рост времени в 14{,}7 раза при переходе от $10^5$ к $10^6$ немного превышает ожидаемые 10 раз. Вероятная причина в том,
что при $n = 10^5$ таблица (около 800~КБ) ещё помещается в кэш процессора, а при $n = 10^6$ и больше выходит за
его пределы, и обращения к памяти становятся медленнее. Это влияет на константу, но не на асимптотику: при дальнейшем
росте $n$ в 10 раз время растёт в 10{,}5 раза.

\pagebreak